\paragraph{Other datasets: MNIST, SVHN and CIFAR-10.}\label{sec:datasets}
Table~\ref{tab:other} repeats the two basic concept-shift scenarios on MNIST, SVHN and CIFAR-10 with the same hyperparameters, and adds DisCo-T30, which uses 30 instead of 10 warm-up rounds.

\paragraph{MNIST.} DisCo-CFL recovers the groups (ARI $1.00$ and $0.99$), never merges conflicting clients, and matches the Oracle (98.1\% and 98.0\%). FL+HC, given $K$, is as good under rotations (ARI $1.00$) but merges 9.4\% of the conflicting pairs under label swaps; all other baselines are 1--14 points below in accuracy and 6--33 points in concept accuracy.

\paragraph{SVHN and CIFAR-10: the reference model is not informative.} On the colour datasets DisCo-CFL with the default warm-up \emph{fails}. On SVHN it returns a single cluster in every run, so it reduces to FedAvg (69.3\% and 73.3\%), whereas the Oracle reaches 86.3\% and 86.5\%: a missed gain of 13--17 points. On CIFAR-10 it forms 1--3 clusters (ARI $0.01$ and $0.28$) and is 4--6 points below the Oracle. The diagnostics locate the cause. After 10 rounds of FedAvg the small CNN has barely started to learn these datasets (on SVHN label swaps, seed 0, the training loss has fallen only from $2.30$ to $2.27$), so the class-conditional gradients of the groups are almost identical: the estimated separation is $\hat\Delta=0.01$ on SVHN and $0.05$ on CIFAR-10, far below the value $1-\tau=0.1$ that Theorem~\ref{thm:rec} requires (Table~\ref{tab:diag}). The method then merges what it cannot distinguish, as designed, but here the groups do need different models. This failure is new in the revision: with the implementation error of the original experiments, clients of a cluster trained sequentially and the 10-round warm-up model was much further trained.

A longer warm-up helps where the model then separates the groups. With 30 rounds (DisCo-T30), $\hat\Delta$ rises to $0.17$ on SVHN label swaps, and DisCo-T30 recovers the groups (ARI $0.81$, 6.6\% conflicting merges) and reaches 85.2\% accuracy, within 1.3 points of the Oracle and above every baseline. On SVHN rotations ($\hat\Delta=0.05$) and on CIFAR-10 ($\hat\Delta=0.08$--$0.09$) 30 rounds are not enough: ARI stays at $0.10$--$0.37$ and accuracy improves by at most $4$ points over the default. No baseline does well there either: the best non-oracle accuracy is 79.2\% on SVHN rotations (PACFL, ARI $0.36$, given $K$) and 66.2\% on CIFAR-10 rotations (PACFL), and all clustering baselines have ARI at most $0.38$.

These results qualify the claim of the paper. DisCo-CFL recovers concept groups when the reference model separates them, and the separation diagnostic tells, after the fact, whether it did. It does not tell, before clustering and without labels, how long the warm-up must be. A rule that stops the warm-up when the training loss reaches a plateau, which we tried on one seed, stops too early on SVHN, because the loss is flat \emph{before} learning starts; we therefore do not report it as part of the method. Choosing the reference model is the main open problem that these experiments expose (Section~\ref{sec:discussion}).

\begin{table}[!htbp]
\centering
\caption{MNIST, SVHN and CIFAR-10 (40 clients, 5 seeds; columns as in Table~\ref{tab:main}). DisCo-T30: DisCo-CFL with 30 warm-up rounds.}\label{tab:other}
{\scriptsize
\begin{tabular}{l r c c c c c c}
\toprule
Method & n & Acc (\%) & Concept (\%) & Worst-grp (\%) & ARI & K & Conflict (\%) \\
\midrule
\multicolumn{8}{l}{\textit{MNIST rotation}} \\
DisCo & 1 & 97.7 & 95.5 & 97.2 & 1.00 & 4.0 & 0.0 \\
FedAvg & 1 & 92.0 & 83.6 & 91.4 & 0.00 & 1.0 & 100.0 \\
Local & 1 & 95.3 & 64.8 & 94.2 & 0.00 & 40.0 & 0.0 \\
IFCA & 1 & 97.8 & 94.9 & 96.7 & 0.93 & 4.0 & 1.7 \\
MTCFL & 1 & 96.6 & 82.1 & 96.0 & 0.11 & 26.0 & 0.0 \\
FedCCFA & 5 & 95.1 $\pm$ 0.4 & 86.1 $\pm$ 1.4 & 94.3 $\pm$ 0.3 & 0.00 $\pm$ 0.00 & 1.6 & 96.0 \\
FedDAG-C & 1 & 97.9 & 95.3 & 97.4 & 1.00 & 4.0 & 0.0 \\
\addlinespace
\multicolumn{8}{l}{\textit{MNIST label swap}} \\
DisCo & 1 & 97.8 & 95.5 & 97.3 & 1.00 & 4.0 & 0.0 \\
FedAvg & 1 & 81.7 & 63.2 & 69.3 & 0.00 & 1.0 & 100.0 \\
FeSEM & 1 & 95.8 & 86.9 & 94.6 & 0.53 & 4.0 & 14.0 \\
FedCCFA & 5 & 96.9 $\pm$ 0.4 & 84.7 $\pm$ 2.6 & 96.2 $\pm$ 0.4 & 0.02 $\pm$ 0.03 & 3.2 & 79.3 \\
\addlinespace
\multicolumn{8}{l}{\textit{SVHN rotation}} \\
FedCCFA & 5 & 71.3 $\pm$ 5.7 & 43.9 $\pm$ 4.9 & 68.4 $\pm$ 7.2 & 0.00 $\pm$ 0.00 & 1.0 & 100.0 \\
\addlinespace
\multicolumn{8}{l}{\textit{SVHN label swap}} \\
Local & 1 & 76.3 & 39.5 & 75.3 & 0.00 & 40.0 & 0.0 \\
MTCFL & 1 & 76.3 & 42.3 & 72.7 & $-$0.03 & 23.0 & 2.7 \\
FedCCFA & 5 & 83.4 $\pm$ 1.2 & 56.9 $\pm$ 1.4 & 81.2 $\pm$ 2.6 & $-$0.00 $\pm$ 0.00 & 1.4 & 97.9 \\
\addlinespace
\multicolumn{8}{l}{\textit{CIFAR-10 rotation}} \\
FeSEM & 1 & 62.2 & 35.3 & 59.5 & 0.07 & 4.0 & 22.7 \\
FedCCFA & 5 & 61.3 $\pm$ 4.1 & 29.2 $\pm$ 0.5 & 57.4 $\pm$ 6.0 & 0.00 $\pm$ 0.00 & 1.4 & 98.0 \\
\addlinespace
\multicolumn{8}{l}{\textit{CIFAR-10 label swap}} \\
FedCCFA & 5 & 64.1 $\pm$ 3.8 & 30.2 $\pm$ 0.8 & 61.2 $\pm$ 5.2 & 0.00 $\pm$ 0.00 & 1.6 & 96.9 \\
\bottomrule
\end{tabular}
}

\end{table}
